\documentclass[10pt,aspectratio=169]{beamer}

% =========================================================
% Academy-style Beamer template (clean + formal)
% Course: 人工智能与数学软件
% =========================================================

% ---------- Fonts & Language ----------
\usepackage[UTF8]{ctex} % Chinese support (XeLaTeX)
\usepackage{fontspec}
\setmainfont{TeX Gyre Heros Cn} % serif-like academic look (safe default)
\setsansfont{WenQuanYi Zen Hei Mono}
\setmonofont{WenQuanYi Zen Hei Mono}

% ---------- Math ----------
\usepackage{amsmath,amssymb,amsthm,mathtools}

% ---------- Figures / Tables ----------
\usepackage{graphicx}
\usepackage{booktabs}

% ---------- Links ----------
\usepackage{hyperref}
\hypersetup{
  colorlinks=true,
  linkcolor=.,
  urlcolor=.,
  citecolor=.
}

% ---------- Beamer Theme (Academy vibe) ----------
\usetheme{Madrid}
\usecolortheme{dolphin}
\usefonttheme{professionalfonts}

% ---------- Tight, clean layout ----------
\setbeamersize{text margin left=10mm,text margin right=10mm}
\setlength{\parskip}{4pt}

% ---------- Footer / Header ----------
\setbeamertemplate{navigation symbols}{}

% Footline: short title | section | slide number
\setbeamertemplate{footline}{
  \leavevmode%
  \hbox{%
    \begin{beamercolorbox}[wd=.38\paperwidth,ht=2.4ex,dp=1.0ex,leftskip=1.2ex]{author in head/foot}
      \usebeamerfont{author in head/foot}\insertshorttitle
    \end{beamercolorbox}%
    \begin{beamercolorbox}[wd=.44\paperwidth,ht=2.4ex,dp=1.0ex,center]{title in head/foot}
      \usebeamerfont{title in head/foot}\insertsectionhead
    \end{beamercolorbox}%
    \begin{beamercolorbox}[wd=.18\paperwidth,ht=2.4ex,dp=1.0ex,rightskip=1.2ex,right]{date in head/foot}
      \usebeamerfont{date in head/foot}\insertframenumber/\inserttotalframenumber
    \end{beamercolorbox}%
  }%
  \vskip0pt%
}

% Optional: subtle headline with section navigation
\setbeamertemplate{headline}{
  \begin{beamercolorbox}[ht=2.2ex,dp=0.8ex,leftskip=1.2ex,rightskip=1.2ex]{section in head/foot}
    \usebeamerfont{section in head/foot}\insertsectionnavigationhorizontal{\paperwidth}{}{}
  \end{beamercolorbox}
}

% ---------- Block style ----------
\setbeamertemplate{blocks}[rounded][shadow=true]

% ---------- Title info ----------
\title[人工智能与数学软件]{\Large 人工智能与数学软件\\\small 用计算机工具（尤其AI）学习数学}
\subtitle{第1周：课程内容工具准备}
\author{王何宇 \\ wangheyu@zju.edu.cn, 13456940632}
\institute{浙江大学数学科学学院}
\date{2026年3月3日}


% =========================================================
\begin{document}

% =========================================================
\begin{frame}[plain]
  \titlepage
\end{frame}

\section{课程简介}
% =========================================================
\begin{frame}{课程简介}
\begin{itemize}
  \item 课程名称：《人工智能与数学软件》
  \item 核心内容：必须掌握的\textbf{技能}、交流使用 AI 的\textbf{经验}和使用 AI 对数学进行\textbf{探索}
  \item 课程组织形式（每周）：2小时课堂讲授 + 2小时讨论/实验
\end{itemize}
\begin{block}{考核方式}
\begin{itemize}
  \item \textbf{无期末考试}
  \item 课程报告（期末提交）：一份结构化、可复现的学习/实验报告
  \item 平时环节：作业、小测（线上/线下）、（期中）项目作业、课堂参与
\end{itemize}
\end{block}
\end{frame}

\section{人工智能对数学的影响}

\begin{frame}{人工智能的非线性快速发展}
\begin{itemize}
    \item GPT 3.5 第一次和公众接触是 2022 年 11 月 30 日
    \item 此后各种大语言模型不但越来越丰富，而且更新周期越来越短
    \item 数学方面：从 1+1 都做不对，到目前数学院硕士研究生水平，\newline
    个别领域达到博士
    \item 那么数学家会不会被取代呢？
    \item<2-> 未来不会使用 AI 的数学家会非常罕见（尽管也许不是完全没有）
\end{itemize}
\end{frame}

\begin{frame}{准备好你的 AI 工具}
    \begin{itemize}
        \item<1-> 免费工具：deepseek, kimi，qianwen，AI 校园(chat.zju.edu.cn)
        \item<2-> 收费工具：askmanyai.cn
        \item<3-> 注意事项：仔细看一下用户协定（换个AI看），明确知识权归属。不要上传涉密材料。
        \item<4-> 暂时掌握网页端使用就可以。
    \end{itemize}

\end{frame}

\begin{frame}{一个完全受自己控制的计算环境}
\begin{itemize}
    \item 推荐：Windows + wsl + debian
    \item 必要的工作环境：python(anaconda) + xelatex(texlive-full) + git
    \item 介绍 overleaf，但本地 latex 不可少
    \item 需要一个编译器，推荐 vs code，装在主机而不是虚拟机中
    \item 编程推荐 python 科学环境：Anaconda 或 Miniconda
\end{itemize}
\end{frame}

\end{document}